\scanpage{55}
% manuscript line 1
que, pour une base orthonormale convenable dans $E$,
$h$ se met sous forme d'une matrice diagonale \,;
\deleted{(\unsure{ses}} \inserted{alors} les \deleted{\unsure{termes}} diagonaux \deleted{sont} non nuls sont
ses valeurs propres \inserted{non nulles $(\rho_i)$} de $h$. On sait alors que $h$ est
de H.-S. si et seulement si cette matrice est
de carré sommable, i.e. si $(\sum\lvert\rho_i\rvert^2)^{1/2}<+\infty$, et
qu'on a alors $\|h\|_2=(\sum\rho_i^2)^{1/2}$. Le th. 2 est par là
complètement démontré.

% manuscript line 9
Les relations entre opérateurs de H.-S. et opérateurs
de Fredholm sont exprimées dans le

% manuscript line 11
\uline{Th. 3} Soient $E$ et $F$ deux espaces de Hilbert. L'application
linéaire canonique de $\bar E\widehat\otimes F$ dans $L(E,F)$ est biunivoque,
\deleted{et} Pour qu'une application linéaire $u$ de $E$ dans $F$
soit application de Fredholm, il faut et il suffit
que $h=\sqrt{u^*u}$ le soit. \deleted{Pour} Pour ceci, il faut et il suffit que
\deleted{et si les $(\lambda_i)$ sont les} $h$ soit compact et la \deleted{\unsure{suite}} $(\rho_i)$ de ces valeurs propres
\deleted{et si les valeurs propres de $h$, \unsure{alors il est} il suffit
que $\sum\lambda_i<+\infty$, et alors on a} \deleted{$\sum\lambda_i\unsure{\geq}0$}
sommable, et alors on a

% manuscript line 20
\[
\|u\|_1=\|h\|_1=\sum\rho_i \tag*{(20)}
\]

% manuscript line 21
On peut aussi y remplacer $h$ par $\sqrt{uu^*}$.
\deleted{Cela \unsure{résulte}}

% manuscript line 22
% margin note written sideways
\margin{La validité des 2 premières \deleted{\unsure{parties}} assertions du théorème a déjà été signalée à la fin du \deleted{au} §2, N\textsuperscript{o} 1, (\uline{petites corrections})}

% manuscript line 23
Que $u$ et $h$ soient simultanément des opérateurs
de Fredholm, et $\|u\|_1=\|h\|_1$, résulte du fait
qu'on a $u=Uh$, $h=Vu$, avec $\|U\|\leq1$, $\|V\|\leq1$.
\deleted{Si $h$ est \unsure{de}} \deleted{Si $h$ est tel \unsure{que} qu'il en est ainsi, et si on} Alors, $h$ sera un opérateur compact, donc
\deleted{$h=\sum\lambda_i\bar a_i\otimes a_i$} de la forme $\sum\rho_i\,\bar a_i\overset{' }{\otimes}a_i$, \deleted{où}
les $(a_i)$ forment un système orthonormal \inserted{dans $E$} et où les
$\rho_i$ sont les valeurs propres de $h$. Soit $p_n$ la
projection de $E$ sur l'espace engendré par les
$a_1,\ldots,a_n$, on aura $\sum_1^n\rho_i=\operatorname{Tr}p_nh\leq\|p_nh\|_1\leq\|h\|_1$
d'où $\sum\rho_i\leq\|h\|_1$. Réciproquement, supposons

% manuscript line 26a
% slanted
\margin{Non !}

% manuscript line 34
% archive sheet number: 56
